\documentclass[11pt]{article}

\usepackage[margin=0.9in]{geometry}
\usepackage{amsmath,amssymb,bm}
\usepackage[hidelinks]{hyperref}

\setlength{\parskip}{0.45em}
\setlength{\parindent}{0em}
\pagestyle{plain}

\begin{document}
\begin{center}
{\Large Bilayer square lattice}\\[0.35em]
{\large Phase diagram from a neural quantum state}
\end{center}

\section{Hamiltonian}

The lattice is a bilayer square lattice. The index $i$ runs over the lattice sites of the square lattice, and $\ell=1,2$ is the layer. The fermion $c_{i\ell\sigma}$ carries spin $\sigma\in\{\uparrow,\downarrow\}$.

The spin on layer $\ell$ of lattice site $i$ is
\begin{equation}
\bm S_{i\ell}=\frac12\sum_{\sigma\sigma'}c_{i\ell\sigma}^{\dagger}\,\bm\sigma_{\sigma\sigma'}\,c_{i\ell\sigma'}.
\end{equation}
Nearest-neighbor pairs on one layer are $\langle ij\rangle$. Next-nearest-neighbor pairs, the diagonals of the plaquettes, are $\langle\langle ij\rangle\rangle$. The Hamiltonian is
\begin{equation}
\begin{aligned}
H={}&-t_1\sum_{\langle ij\rangle,\ell,\sigma}\bigl(c_{i\ell\sigma}^{\dagger}c_{j\ell\sigma}+\mathrm{h.c.}\bigr)
-t_2\sum_{\langle\langle ij\rangle\rangle,\ell,\sigma}\bigl(c_{i\ell\sigma}^{\dagger}c_{j\ell\sigma}+\mathrm{h.c.}\bigr)\\
&+J\sum_{i}\bm S_{i1}\cdot\bm S_{i2}.
\end{aligned}
\end{equation}
There is no interlayer hopping. $J>0$ is antiferromagnetic. Set $t_1=1$. This fixes the energy unit. Set $t_2=0.2$.

Each lattice site has two layers and two spin states, four single-particle states in all. The filling $\nu$ is the number of electrons divided by four times the number of lattice sites. Equivalently, $\nu$ is the electron number per layer, per spin, divided by the number of lattice sites. Full filling is $\nu=1$: both layers and both spins are occupied. Half filling is $\nu=1/2$. Quarter filling is $\nu=1/4$. The particles need not be polarized; $\nu$ fixes only the density.

\section{Task}

Solve $H$ by variational Monte Carlo with a neural-network wavefunction. State the lattice size, the boundary conditions, the ansatz, the sampler, the optimizer, and the statistical uncertainty on every reported number.

\subsection{Line cuts}

Compute two cuts through the ground-state phase diagram, each against $J$, from weak coupling through strong coupling. One cut is at half filling, $\nu=1/2$. The other is at quarter filling, $\nu=1/4$. For each point, report the energy per lattice site and the diagnostics that decide the phase: a charge gap, a spin gap, a magnetic order parameter, a superconducting pair correlation, and the momentum distribution. Say whether that point is a Fermi liquid, a gapped state that preserves every symmetry of $H$, a superconductor, or another phase. A label without the diagnostic is not an answer.

On the half-filling cut, identify whether a gapped phase is present.

\subsection{Phase diagram}

Repeat the identification on the plane of $\nu$ and $J$, for $0<\nu\le 1$. The required figure is that phase diagram, with the two line cuts marked on it.

\subsection{Spectra of the gapped phase}

At one point inside the gapped phase found on the half-filling cut, compute three spectra.

\begin{itemize}
\item The charge spectrum measured by scanning tunneling: differential conductance versus bias.
\item The spin excitation spectrum that enters muon spin relaxation or nuclear spin-lattice relaxation, as a function of excitation energy.
\item The resistivity, or the conductivity, as a function of temperature.
\end{itemize}

If a spectrum is not a direct ground-state expectation value, state the approximation that produces it. Give every energy in units of $t_1$.

\end{document}
